\chapter{引言与问题定义}

\section{算例背景}
平板泊肃叶流是最适合做方法验证的内部流动之一。几何简单、解析解明确、稳态流场平滑，既能检查数值离散对粘性层流的恢复能力，也能把误差来源压缩到核函数、边界处理、时间推进和粒子分布控制这些核心模块。对 SPH 这类无网格方法而言，它还是一个很典型的基准问题：横向存在固壁约束，纵向存在周期延拓，流动由恒定体积力驱动，最终稳态速度剖面应恢复为抛物线形。

当前仓库实现的是二维弱可压 SPH 平板泊肃叶流。最近一次提交把旧的多分支壁面策略收敛为“单层 shell 壁面 + 唯一 no-slip”实现，同时把主程序进一步整理为单个 \code{SPH\_Poiseuille.m} 壳层文件：配置、初始化、主循环框架、restart、监控和后处理都仍保留在 MATLAB 侧，但长段时间推进与壁面剪应力统计已经通过高层 MEX 接口下沉到 C 层。程序除了输出稳态速度剖面与速度场，还记录通道中截面的速度演化，并在运行期打印上下壁面切应力估计值，方便判断数值解是否稳定地向解析解收敛。

\section{报告目标}
本报告不把仓库视为一个黑盒算例，而是把它拆成五个彼此衔接的层次来说明。第一层是物理模型，即控制方程、体积力与解析解之间的关系；第二层是 SPH 离散，即核函数、粒子对、密度重初始化、粘性项和压力项；第三层是算法调度，即统一 CFL 时间步、两阶段积分以及输出节拍如何组织；第四层是边界与稳定性控制，即周期边界、单层 shell 壁面、无滑移接触和防穿透；第五层是实现架构，即 MATLAB 脚本与 C/MEX 文件之间如何分工，以及现有结果图如何生成。

面向项目维护，这份报告还有一个额外任务：将“文档说法”和“代码真相”对齐。最近一次提交之后，\code{README.md}、历史设计文档和旧报告里的若干表述已经过时，例如 \code{bc\_mode}、PI 控制、Neumann 风格壁面分支、双层时间步进以及独立的分辨率扫描脚本。报告会把这些地方统一改写为当前仓库真实存在的实现路径，避免后续扩展工作基于过时描述做错误判断。

\section{整体计算流程}
从流程上看，程序仍然可以分成七个阶段。第一阶段自动检查并编译 \code{sph\_neighbor\_search\_mex} 和运行时物理 MEX \code{sph\_physics\_shell\_mex}。第二阶段解析配置文件，计算派生量，例如 $g$、$U_{max}$、$h$ 和 shell 壁面厚度。第三阶段按均匀网格生成流体粒子，并调用 \code{build\_shell\_wall\_particles.m} 在上下壁外各布置一层中面壁面粒子。第四阶段检查 restart 文件并在签名匹配时恢复状态。第五阶段建立初始粒子对，完成密度重初始化与核梯度修正，并把运行数据收拢为 \code{cfg/state/geom/neighbor/monitor}。第六阶段进入主循环，在每个统一时间步中先通过 \code{advance\_one\_step\_mex} 调用 C 层门面 \code{advance\_shell\_step} 完成单步推进，再通过 \code{wall\_shear\_monitor\_mex} 调用 \code{wall\_shear\_monitor} 完成上下壁面剪应力统计，随后执行周期包裹、壁面约束、排序与邻居重建。第七阶段进行后处理，计算速度剖面、L2 误差并输出两张结果图。

这个七阶段结构有两个直接好处。其一，源码分层清楚，MATLAB 主文件仍保持了足够强的可读性。其二，MEX 的职责边界也更加明确：邻居搜索只做几何关系构造，物理 MEX 同时提供底层离散算子和高层推进门面，而时间推进策略、输出节拍和结果管理仍留在主脚本层。

\section{报告组织}
后续章节按“从物理到代码”的顺序展开。第 \ref{chap:physical} 章说明物理模型和解析关系，第 \ref{chap:sph} 章说明 SPH 离散公式，第 \ref{chap:time} 章说明统一 CFL 时间推进与两阶段积分，第 \ref{chap:boundary} 章说明周期边界、单层 shell 壁面与防穿透，第 \ref{chap:code} 章给出文件级和函数级映射，第 \ref{chap:results} 章解释现有数值结果与运行期验证指标，第 \ref{chap:discussion} 章总结当前实现的优点、局限与后续建议。附录则给出关键代码索引和推荐运行命令，方便读者回到仓库直接定位。
